\section{Post-Deployment Diagnostic Simulation \diagmark}
\label{sec:diagnostic}

\subsection{Why this is a separate kind of evidence}

Everything in this section was run after the physical campaign, by people
who already knew how it had gone. It is not a prediction, it was not
pre-registered, and no result here may be read as confirming the frozen
study. We keep the distinction typographic (\diagmark\ against
\frozenmark) rather than relying on the reader to remember which section
they are in.

The motivation is that the physical campaign cannot answer its own most
interesting question. Seven changes were made to the stack in a single
morning, all of them together, with no repetitions of the intermediate
states, because there was one day of pool time. Which of them carried the
effect is unanswerable from eight runs. Simulation can do what the pool
cannot: apply the changes one at a time, with repetitions, holding
everything else at its frozen value.

\subsection{A self-contained baseline, and why the frozen numbers are not
used as one}

The natural design would compare each rung against the 80 frozen runs.
We do not, because the frozen campaign is not reproducible. Two causes
are identified in Section~\ref{sec:limitations}: the observation relay's
random seed is discarded when it is zero, which is the value the campaign
used; and the frozen campaign's command loop ran at roughly half its
intended rate because that loop is wall-clock driven and twenty runs were
executed back to back on a loaded machine. Re-running the frozen
configuration on an idle machine gives minimum clearances spanning
\SIrange{0.29}{3.25}{\metre} where the frozen campaign spanned
\SIrange{0.50}{0.65}{\metre}.

The ablation is therefore self-contained. Its first rung is the frozen
\emph{configuration} run in this session, on this machine, alongside
every other rung, and every comparison is rung-against-rung. Ten
repetitions per rung are used rather than five, because the baseline's
own spread demands it, and results are reported as medians with
interquartile ranges together with the fraction of runs that manoeuvred
at all -- a bimodal outcome is not honestly summarised by a mean.

\subsection{The ladder}

Each rung adds one documented deployment change to the frozen
configuration; the final rung is all of them together, which is the
configuration that flew.

\begin{table}[t]
\centering
\caption{The diagnostic ablation \diagmark. All runs use the committed
planner, the \textsf{K0} geometry and the S3 calibration level, with ten
repetitions per rung. Rungs A1--A4 are changes made in the pool; A5 is a
defect found in the frozen study during the final audit; A6 is a
constraint nobody noticed on the day.}
\label{tab:ablation-design}
\begin{tabular}{lll}
\toprule
rung & change & origin \\
\midrule
A0 & frozen configuration & baseline \\
A1 & qualifier warm-up $20\!\to\!2$, confirm $3\!\to\!1$ & forced in pool \\
A2 & risk thresholds $0.55/0.30 \to 0.30/0.15$ & forced in pool \\
A3 & avoidance hold $1.0 \to \SI{4.0}{\second}$ & forced in pool \\
A4 & engagement range $1.5 \to \SI{1.8}{\metre}$ & as actually flown \\
A5 & relay field of view $60 \to \SI{90}{\degree}$ & audit defect \\
A6 & clearing offset $2.5 \to \SI{0.6}{\metre}$, run-past
     $4.0 \to \SI{1.0}{\metre}$ & basin constraint \\
A7 & all of the above & the flown configuration \\
\bottomrule
\end{tabular}
\end{table}

A4 uses \SI{1.8}{\metre} and not the \SI{5.0}{\metre} recorded in the
session file, for the reason given in Section~\ref{sec:deployment}: every
run records \SI{1.8}{\metre}, and the session file is wrong.

A6 is not a change anyone made. It is included because the frozen planner
asks for a manoeuvre the basin cannot contain, and
Section~\ref{sec:simreal} shows from the command stream that no physical
run ever left the strafe phase. A6 asks what the same planner does when
given a manoeuvre that fits, and is the rung we expect to matter most --
a prediction we record here, before reporting the outcome, because the
rest of this section was written with the physical results already known.

\subsection{Results}

Table~\ref{tab:ablation} reports the ablation. Each rung is ten
runs of the committed planner in \textsf{K0} at S3, and every
comparison is against A0 in the same session.


\begin{table}[t]
\centering
\caption{Deployment-gap ablation \diagmark, ten runs per rung, medians
with interquartile range in brackets. \emph{go-around} counts runs that
reached the planner's second manoeuvre phase, the observable used on the
physical runs in Section~\ref{sec:simreal}; no physical run reached it.}
\label{tab:ablation}
\setlength{\tabcolsep}{4pt}
\begin{tabular}{llrrrrr}
\toprule
rung & change & clearance (\si{\metre}) & lateral (\si{\metre})
& manoeuvred & go-around & collisions \\
\midrule
A0 & frozen configuration & 1.29 \tiny[0.42--2.66] & 1.67 \tiny[1.35--3.96] & 8/10 & 10/10 & 2 \\
A1 & qualifier relaxed & 2.22 \tiny[1.33--3.12] & 2.76 \tiny[2.21--4.47] & 9/10 & 10/10 & 2 \\
A2 & risk thresholds lowered & 1.90 \tiny[0.73--2.23] & 1.98 \tiny[1.29--3.31] & 9/10 & 10/10 & 2 \\
A3 & avoidance hold lengthened & 2.43 \tiny[1.79--3.15] & 1.73 \tiny[1.28--4.23] & 9/10 & 10/10 & 0 \\
A4 & engagement range as flown & 1.48 \tiny[0.30--2.62] & 1.52 \tiny[1.20--2.06] & 10/10 & 10/10 & 3 \\
A5 & relay field of view corrected & 2.91 \tiny[2.41--3.23] & 1.56 \tiny[0.03--2.73] & 7/10 & 10/10 & 0 \\
A6 & pool-feasible manoeuvre & 0.30 \tiny[0.28--2.95] & 0.47 \tiny[0.00--0.82] & 7/10 & 10/10 & 4 \\
A7 & the flown configuration & 3.13 \tiny[2.84--3.14] & 2.33 \tiny[0.85--6.02] & 10/10 & 6/10 & 2 \\
\bottomrule
\end{tabular}
\end{table}

Four results stand out, and the first two were not anticipated.

\emph{Our own frozen-campaign defect was the single largest effect on
safety.} Correcting the relay's field of view (A5) raises median
clearance from \SI{1.29}{\metre} to \SI{2.91}{\metre} and removes both
collisions. The defect made the relay believe the obstacle roughly
\num{1.5} times further away than the planner did, so the measured
detection-probability curve was applied in the wrong band and
observations were suppressed where they mattered most. The frozen
predictions are not wrong about the ladder they describe; they are a
description of a simulator with a bug in it, and the bug cost more
clearance than any calibration level added.

\emph{The manoeuvre that fits the basin is not safe.} A6 gives the
planner a clearing offset and run-past margin that a
\SIrange{2.7}{3.0}{\metre} basin can contain. Median clearance collapses
to \SI{0.30}{\metre} and four runs of ten end in collision -- the worst
of any rung. The committed planner's safety margin comes from the width
of its detour, and there is no version of that detour that both fits this
pool and keeps the vehicle away from the anchor. This is an architectural
result, not a tuning one: the physical campaign was not a fair test of
this planner, because the basin cannot accommodate the manoeuvre the
planner is built around.

\emph{The changes forced in the pool were individually benign.}
Relaxing the qualifier (A1), lowering the risk thresholds (A2) and
lengthening the avoidance hold (A3) all raise median clearance relative
to A0, and A3 removes collisions entirely. Engaging at the distance
actually flown (A4) is the only forced change that degrades safety, and
mildly. The morning of adaptation was not, in simulation, the source of
the deployment gap.

\emph{The simulated world has no walls, so it cannot reproduce what
reality did.} Every rung reaches the go-around phase in ten runs of ten,
except the fully deployed configuration at six of ten. No physical run
reached it at all (Section~\ref{sec:simreal}). The reason is not a
parameter: the simulated scenario is open water with an obstacle in it,
so a \SI{2.5}{\metre} strafe always completes, while in the pool it
cannot. The basin boundary is absent from the simulator and from all
three calibration levels, and it is the structural mismatch that most
shaped the physical runs. We report it as the clearest instance of this
paper's general point -- that the mismatches which mattered were not the
ones we measured.

